\documentclass[11pt]{article}

\usepackage[T1]{fontenc}
\usepackage[margin=1in]{geometry}
\usepackage{amsmath,amssymb,amsthm,mathtools}
\usepackage{hyperref}
\usepackage{enumitem}

\newtheorem{theorem}{Theorem}
\newtheorem{lemma}[theorem]{Lemma}
\newtheorem{proposition}[theorem]{Proposition}
\newtheorem{corollary}[theorem]{Corollary}
\newtheorem{remark}[theorem]{Remark}

\newcommand{\Pset}{\mathbb P}
\newcommand{\Z}{\mathbb Z}
\newcommand{\R}{\mathcal R}
\newcommand{\B}{\mathcal B}
\newcommand{\C}{\mathcal C}
\newcommand{\Hcal}{\mathcal H}
\newcommand{\eps}{\varepsilon}
\newcommand{\one}{\mathbf 1}

\title{A proof of Erd\H{o}s Problem 689}
\author{Working note}
\date{April 26, 2026}

\begin{document}
\maketitle

\begin{abstract}
We prove that for all sufficiently large $n$ one can choose one residue class
$a_p\pmod p$ for each prime $p\le n$ so that every integer $m\in[1,n]$ lies in
at least two of the chosen classes.  The proof starts with the parity-first
assignment $a_2\equiv1\pmod2$, leaves $3$ in the zero class, switches a fixed
large set of small primes $S\subset\{7,11,13,\ldots\}$ to nonzero classes, and
then uses robust cleanup primes $P>n/5$.  The cleanup is reduced to a
finite-core prime-difference matching problem.  A half-residue transport kernel
produces a fractional matching with fixed slack; the finite-complexity
Green--Tao--Ziegler theorem for affine-linear forms in primes gives the
required first and second weighted moments; and Kahn's fractional
Frankl--R\"odl--Pippenger theorem rounds the fractional matching to an integral
matching.  The exact form of Kahn's theorem used here is Theorem 1.5 of
\cite{Kahn1996}.
\end{abstract}

\section{The parity-first residual problem}

For a prime $p\le n$, write $a_p\pmod p$ for the residue class chosen for $p$.
We begin with
\[
  a_2\equiv1\pmod2,
\]
leave all odd primes in the zero class, and later switch selected odd primes to
nonzero classes.  The baseline coverage is
\[
  C_0(m)=\one_{m\text{ odd}}+\#\{p\text{ odd prime}:p\mid m\}.
\]
If $R$ is the set of switched odd primes, $b_p$ is the new residue for $p\in R$,
\[
  L_R(m)=\#\{p\in R:p\mid m\},\qquad
  G_R(m)=\#\{p\in R:m\equiv b_p\pmod p\},
\]
then the exact requirement is
\[
  G_R(m)\ge \max(0,2-C_0(m)+L_R(m)).       \tag{1}
\]
Thus switching a prime creates debt only at multiples whose previous coverage
used that prime's zero class.

We keep $3$ at $0\pmod3$.  Choose a fixed finite set
\[
  S\subset\{7,11,13,\ldots\}
\]
and switch each $s\in S$ to a nonzero residue $b_s\pmod s$.  Put
\[
  H_S(x)=\#\{s\in S:x\equiv b_s\pmod s\}.
\]
After these fixed switches, the main residual tokens are even integers
\[
  x=2^k u q,
  \qquad k\ge1,                                      \tag{2}
\]
where $u$ is odd and $S$-smooth, $q\notin S$ is an odd prime occurring to
exponent one, and $H_S(x)=0$.  Prime-power exceptions, pure
$\{2\}\cup S$-smooth terms, $m=1$, and other boundary cases contribute
$o(n/\log n)$ residual tokens.

Let $A_S(n)$ denote the main residual token set.

\begin{lemma}[Residual asymptotic and exceptional tokens]\label{lem:residual}
For fixed $S$,
\[
  |A_S(n)|=(1+o(1)){n\over\log n}.                  \tag{3}
\]
Moreover the number of residual tokens not included in $A_S(n)$ is
\[
  O_S\!\left(\sqrt n+(\log n)^{|S|+1}\right)
  =o(n/\log n).
\]
Finally, if a finite coefficient core omits coefficient weight at most
$\eps$ in the sense of the proof below, then the number of omitted main
residual targets is $O_S(\eps n/\log n)+o(n/\log n)$.
\end{lemma}

\begin{proof}
Let $\mathcal D_S$ be the odd $S$-smooth numbers.  For
$d\in\mathcal D_S$ put
\[
  \Theta_S(d)=\prod_{\substack{s\in S\\s\nmid d}}{s-2\over s-1}.
\]
For fixed $d$ and $k\ge1$, the residuals $2^kdq\le n$ with $q\notin S$ prime
are counted in a fixed union of reduced residue classes modulo $W$, hence
fixed-modulus PNT gives
\[
  \#\{q:2^kdq\le n,\ H_S(2^kdq)=0\}
  =
  \Theta_S(d){n\over 2^kd\log n}
  +o_{d,k}(n/\log n).                                \tag{3a}
\]
The coefficient sum is
\[
  \sum_{d\in\mathcal D_S}{\Theta_S(d)\over d}
  =
  \prod_{s\in S}\left({s-2\over s-1}+\sum_{e\ge1}{1\over s^e}\right)=1,
  \qquad
  \sum_{k\ge1}2^{-k}=1.                              \tag{3b}
\]
To justify summing the asymptotics, choose a finite set
$\mathcal K\subset \mathcal D_S\times\mathbb N$ so that
\[
  \sum_{(d,k)\notin\mathcal K}{\Theta_S(d)\over 2^kd}<\eps.
\]
On the finite set $\mathcal K$, (3a) may be summed directly.  On the tail with
$2^kd\le n^{1/2}$, Brun--Titchmarsh or the standard fixed-modulus upper bound
$\pi(x;W,a)\ll_S x/\log x$ gives $O_S(\eps n/\log n)$.  On the tail with
$2^kd>n^{1/2}$, the trivial bound by $\pi(n/(2^kd))\le n^{1/2}$ and the
estimate
\[
  \#\{(d,k):d\in\mathcal D_S,\ 2^kd\le n\}
  \ll_S(\log n)^{|S|+1}
\]
give $o(n/\log n)$.  Letting $\eps\to0$ proves (3).

It remains to bound exceptional tokens.  Odd $S$-smooth integers and even
$\{2\}\cup S$-smooth integers contribute
$O_S((\log n)^{|S|+1})$ tokens.  Integers of the form
$2^kdq^a\le n$ with $a\ge2$, $d\in\mathcal D_S$, and $q\notin S$ prime
contribute at most
\[
  \sum_{d,k} \#\{q^a\le n/(2^kd):a\ge2\}
  \ll_S
  \sqrt n\sum_{d\in\mathcal D_S}{1\over\sqrt d}
  \sum_{k\ge1}2^{-k/2}
  \ll_S\sqrt n.
\]
This proves the exceptional-token bound.  The finite-core tail estimate is
the same tail argument with $\mathcal K$ chosen to be the coefficient core.
\end{proof}
Define
\[
  A_1(n)=\{x\in A_S(n):v_2(x)=1\},\qquad
  A_2(n)=\{x\in A_S(n):v_2(x)\ge2\}.
\]
Then
\[
  |A_1(n)|=\left(\frac12+o(1)\right)|A_S(n)|,
  \qquad
  |A_2(n)|=\left(\frac12+o(1)\right)|A_S(n)|.        \tag{4}
\]

\section{Robust cleanup primes}

Let
\[
  W=\prod_{s\in S}s.
\]
For a unit class $\pi\pmod W$, call $\pi$ robust if
\[
  H_S(\pi)\ge1,\qquad H_S(2\pi)\ge2,
  \qquad H_S(4\pi)\ge2.                              \tag{5}
\]
Let $\B\subset(\Z/W\Z)^\times$ be the robust set and set
\[
  \delta_S=\frac{|\B|}{\varphi(W)}.
\]
For $s\in S$, the three local classes
\[
  \pi\equiv b_s,
  \qquad 2\pi\equiv b_s,
  \qquad 4\pi\equiv b_s\pmod s
\]
are distinct because $s\ge7$.  Hence $\delta_S$ depends only on $S$, not on the
particular nonzero residues $b_s$.

\begin{lemma}[Robust density]
There are fixed finite sets $S\subset\{7,11,13,\ldots\}$ for which $\delta_S$ is
arbitrarily close to $1$.
\end{lemma}

\begin{proof}
Let
\[
  A_S^{(0)}=\prod_{s\in S}\left(1-\frac1{s-1}\right),
  \qquad
  B_S=\sum_{s\in S}\frac1{s-2}.
\]
If $X_1,X_2,X_4$ count hits to $\pi,2\pi,4\pi$, then
\[
  \delta_S=\mathbb P(X_1\ge1,X_2\ge2,X_4\ge2).
\]
Now
\[
  \mathbb P(X_1=0)=A_S^{(0)},
\]
and
\[
  \mathbb P(X_2<2)=A_S^{(0)}(1+B_S),
\]
with the same estimate for $X_4$.  Thus
\[
  \delta_S\ge1-A_S^{(0)}(3+2B_S).                    \tag{6}
\]
Taking $S(y)=\{p:7\le p\le y,\ p\text{ prime}\}$, Mertens' theorem and the
prime number theorem give $A_{S(y)}^{(0)}\asymp1/\log y$ and
$B_{S(y)}=O(\log\log y)$, so the right side of (6) tends to $1$.
\end{proof}

Set
\[
  \beta_*={1\over2}\left(1-{3\over5}e^{-2}\right)=0.459399\ldots,
\]
\[
  \delta_*={1\over \beta_*+3/5}=0.943931\ldots.
\]
Choose $S$ with $\delta_S>\delta_*$, and choose
\[
  \delta_S^{-1}-{3\over5}<\beta<\beta_*.             \tag{7}
\]
Let
\[
  \Delta=(\beta+3/5)\delta_S-1>0.                    \tag{8}
\]
This fixed positive margin absorbs coefficient tails and exceptional tokens.
Define
\[
  \R_\beta(n)=\{P\in(n/5,\beta n]:P\text{ prime},\ P\bmod W\in\B\}
\]
and
\[
  \R(n)=\{P\in(n/5,n]:P\text{ prime},\ P\bmod W\in\B\}.
\]
Fixed-modulus PNT gives
\[
  |\R_\beta(n)|=\left((\beta-1/5)\delta_S+o(1)\right){n\over\log n},
  \tag{9}
\]
\[
  |\R(n)|=\left({4\over5}\delta_S+o(1)\right){n\over\log n}.            \tag{10}
\]

\begin{lemma}[Robust side debt]
Switching any robust prime $P>n/5$ to any nonzero residue creates no unresolved
side debt.
\end{lemma}

\begin{proof}
The only multiples of $P$ not exceeding $n$ are $P,2P,3P,4P$.  The integer $P$
is odd and has the parity hit plus at least one $S$-hit.  The integers $2P$ and
$4P$ have at least two $S$-hits by (5).  The integer $3P$ has the parity hit and
the unchanged zero class modulo $3$.  There is no multiple $5P\le n$.
\end{proof}

\section{The finite-core hypergraph and the half-residue kernel}

Choose finite coefficient cores
\[
  \mathcal A_1\subset\{a:a\text{ odd and }S\text{-smooth}\},
\]
\[
  \mathcal A_2\subset\{b:b=2^j u,\ j\ge1,
       u\text{ odd and }S\text{-smooth}\}.
\]
Let $X_n\subset A_1(n)$ be the residual targets $x=2aq$ with
$a\in\mathcal A_1$, and let $Y_n\subset A_2(n)$ be the residual targets
$y=2bq'$ with $b\in\mathcal A_2$.  The cores will be chosen so that they capture
fractions $\alpha_X,\alpha_Y$ of the full coefficient mass, uniformly in the
half-residue classes below, and
\[
  \alpha_X,
  \alpha_Y>G(\beta)+\eta                                    \tag{11}
\]
for some fixed $\eta>0$, where
\[
  G(\beta)=\int_{1/5}^{\beta}{dt\over1-2t}
  ={1\over2}\log\left({3/5\over1-2\beta}\right)<1.            \tag{12}
\]
This is possible because $\beta<\beta_*$.
In the explicit notation below this means
\[
  \alpha_X=\min_{A\in\C}2\Xi_A,\qquad
  \alpha_Y=\min_{A\in\C}2H_A.
\]
The quantifiers are chosen in this order: first choose $S,\delta_S,\beta$ with
strict margin in (7), then choose $\eta>0$ with $G(\beta)+\eta<1$, and finally
choose the finite cores so that $\alpha_X,\alpha_Y>G(\beta)+\eta$ and the
omitted coefficient-tail mass is smaller than the cleanup surplus.

Let
\[
  c_s\equiv2^{-1}b_s\pmod s,
\]
and define
\[
  \C=\{A\pmod W:A\not\equiv c_s\pmod s\text{ for every }s\in S\}.
\]
For $x=2aq$ write $A\equiv aq\pmod W$, and for $y=2bq'$ write
$B\equiv bq'\pmod W$.  Then residual membership is $A,B\in\C$.

\begin{lemma}[Half-residue regularity]
For every unit $\pi\pmod W$ and every $\sigma\in\{\pm1\}$,
\[
  \#\{A\in\C:A+\sigma\pi\in\C\}=M:=\prod_{s\in S}(s-2).       \tag{13}
\]
Moreover, the full infinite $S$-smooth coefficient distribution is uniform on
$\C$ on both the $X$- and $Y$-sides, and finite cores can approximate this
uniformly.
\end{lemma}

\begin{proof}
For (13), modulo each $s\in S$ the conditions are
$A\ne c_s$ and $A\ne c_s-\sigma\pi$.  Since $\pi$ is a unit, these are two
distinct forbidden classes, leaving $s-2$ choices.

For the coefficient distribution, fix $s\in S$.  If the $s$-adic exponent of the
coefficient is $0$, then for every nonzero half-residue $A\pmod s$ there is one
reduced class of the prime variable, giving local mass $1/(s-1)$.  If the
exponent is at least $1$, then $A\equiv0\pmod s$ and the coefficient weight is
$\sum_{e\ge1}s^{-e}=1/(s-1)$, with the prime variable unrestricted among
reduced classes.  Thus every allowed local half-residue has the same mass.
Multiplying over $s\in S$ gives uniformity on $\C$.  Truncating exponents to
$e\le E$ changes only the $A\equiv0\pmod s$ local mass by $O(s^{-E})$, so since
$S$ is fixed, large finite cores approximate the uniform distribution uniformly.
The extra $2$-power in the $Y$-coefficients contributes a common factor
$\sum_{1\le j\le J}2^{-j}$ and does not affect half-residue uniformity.
\end{proof}

Define the 3-partite 3-uniform hypergraph
\[
  H_n\subset X_n\times Y_n\times Z_n,
  \qquad Z_n=\R_\beta(n),
\]
by putting an edge $(x,y,P)$ when
\[
  |y-x|=2P.                                                 \tag{14}
\]
The parity split is essential: $X_n$ consists of targets $v_2=1$, while $Y_n$
consists of targets $v_2\ge2$, so $(y-x)/2$ is odd.

The aggregate continuum transport is as follows.  A label is $(t,\pi)$ with
$t\in(1/5,\beta]$ and $\pi\in\B$.  Choose $\sigma=\pm1$ with probability $1/2$,
choose uniformly an $A\in\C$ with $A+\sigma\pi\in\C$, and choose
$u\in(0,1-2t)$ uniformly.  If $\sigma=+1$, send mass to
\[
  X=(u,A),\qquad Y=(u+2t,A+\pi),
\]
and if $\sigma=-1$, send mass to
\[
  X=(u+2t,A),\qquad Y=(u,A-\pi).
\]
This gives exact label load $1$.  The side load at a half-residue class is at
most
\[
  \int_{1/5}^{\beta}{dt\over1-2t}=G(\beta).
\]
For finite cores this is multiplied by $\alpha_X^{-1}$ or $\alpha_Y^{-1}$, hence
by (11) the limiting side loads are at most $1-2\gamma$ for some fixed
$\gamma>0$.

The aggregate transport lifts to typed kernels.  A type is
\[
  \tau=(a,b,\sigma,r,r',\pi),
\]
where $a\in\mathcal A_1$, $b\in\mathcal A_2$,
\[
  q\equiv r\pmod W,
  \qquad q'\equiv r'\pmod W,
\]
\[
  \pi\equiv\sigma(br'-ar)\pmod W,
  \qquad \pi\in\B.
\]
The support polygon is
\[
  D_\tau=\left\{(Q,Q'):
  0<Q\le{1\over2a},\ 0<Q'\le{1\over2b},\;
  {1\over5}<\sigma(bQ'-aQ)\le\beta\right\}.              \tag{15}
\]
Every residue-compatible type is admissible.  Indeed, modulo $2$ the parity
split gives $P=\sigma(bq'-aq)$ odd; modulo every $s\in S$, the prescribed
classes have $q,q'$ and $\pi$ units; and outside $2W$ the three forms
$q,q',\sigma(bq'-aq)$ are three distinct lines, so their zero sets cannot cover
all residue pairs.  Let $\lambda_\tau>0$ be the fixed edge-total constant in the
Green--Tao--Ziegler asymptotic.  The typed kernel $g_\tau$ is the lifted
aggregate density divided by $\lambda_\tau$.  It is bounded because the
coefficient core is finite, $\beta<1/2$, and every $\lambda_\tau$ is fixed and
positive.  The condition $\gcd(a,b)=1$ is automatic:
if some $s\in S$ divided both $a$ and $b$, then $A\equiv B\equiv0\pmod s$, but
admissibility gives $B-A\equiv\sigma\pi\not\equiv0\pmod s$.

We record the typed lift explicitly.  For a half-residue $A\in\C$ set
\[
  \Xi_A=\sum_{\substack{a\in\mathcal A_1,\ r\\ ar\equiv A\ (W)}}{1\over2a},
  \qquad
  H_A=\sum_{\substack{b\in\mathcal A_2,\ r'\\ br'\equiv A\ (W)}}{1\over2b}.
\]
The finite cores are chosen so that
\[
  {1-\eps\over2}\le \Xi_A,H_A\le {1\over2}\qquad(A\in\C),       \tag{15a}
\]
with $\eps>0$ small enough that $1-\eps>G(\beta)+\eta/2$.  Equivalently,
$\alpha_X,\alpha_Y\ge1-\eps$, so this is the uniform version of (11).  Define conditional
weights
\[
  \rho_X^A(a,r)={1/(2a)\over\Xi_A},\qquad
  \rho_Y^B(b,r')={1/(2b)\over H_B}.
\]
For $t\in(1/5,\beta]$ and $A+\sigma\pi\in\C$, let
\[
  K_\sigma(t,\pi;u,A)={1\over2M(1-2t)}
  \quad(0<u<1-2t),
\]
and let it be zero otherwise.  If $(Q,Q')\in D_\tau$, put
\[
  t=\sigma(bQ'-aQ),\qquad
  u_\tau(Q,Q')=
  \begin{cases}
    2aQ,&\sigma=+1,\\
    2bQ',&\sigma=-1.
  \end{cases}
\]
The change of variables $(Q,Q')\mapsto(u_\tau,t)$ has Jacobian $2ab$.
With the GTZ normalization below, the base edge measure of type $\tau$ is
$\lambda_\tau\,dQ\,dQ'$.  The typed kernel is therefore
\[
  g_\tau(Q,Q')
  =
  {2ab\over\lambda_\tau}\,
  K_\sigma(t,\pi;u_\tau(Q,Q'),A)\,
  \rho_X^A(a,r)\rho_Y^B(b,r').                         \tag{15b}
\]
This formula gives the load identities.  For labels, the coarea factor is
$b^{-1}dQ$, and after summing the conditional distributions over all typed
fibres above a fixed half-residue pair $(A,A+\sigma\pi)$ one obtains
\[
  L_Z^{\rm lim}(\pi,t)
  =
  \sum_{\sigma=\pm1}
  \sum_{\substack{A\in\C\\A+\sigma\pi\in\C}}
  \int_0^{1-2t}K_\sigma(t,\pi;u,A)\,du=1.              \tag{15c}
\]
For the $X$-side, writing $z=2aQ$, the same change of variables gives
\[
  L_X^{\rm lim}(a,r,Q)
  =
  {1\over\Xi_A}
  \left[
    {N_+(A)\over2M}J_\beta\!\left({1-z\over2}\right)
    +
    {N_-(A)\over2M}J_\beta\!\left({z\over2}\right)
  \right],                                             \tag{15d}
\]
where
\[
  N_\pm(A)=\#\{\pi\in\B:A\pm\pi\in\C\},\qquad
  J_\beta(u)=\int_{1/5}^{\min(\beta,u)}{dt\over1-2t}
\]
with the integral interpreted as $0$ if $u\le1/5$.  Since $N_\pm(A)\le M$ and
(15a) holds,
\[
  L_X^{\rm lim}
  \le {1\over1-\eps}
  \left[
    J_\beta\!\left({1-z\over2}\right)
    +J_\beta\!\left({z\over2}\right)
  \right]
  \le {G(\beta)\over1-\eps}\le1-2\gamma.
\]
The middle inequality uses $\beta>2/5$, which follows from (7) and
$\delta_S\le1$: if $u+v=1/2$, then
$J_\beta(u)+J_\beta(v)\le G(\beta)$.
The $Y$-side is identical with $\Xi_A$ replaced by $H_B$.  This proves the
finite-core typed-kernel lift and the side slack used below.

\section{GTZ weighted moments}

For an edge $e=(x,y,P)$ of type $\tau$, write
\[
  x=2aq,
  \qquad y=2bq',
  \qquad P=\sigma(bq'-aq),
\]
and assign
\[
  w_e={\log^2 n\over n}g_\tau(q/n,q'/n).                    \tag{16}
\]
The following proposition is the only analytic input.

\begin{proposition}[Weighted GTZ moment proposition]\label{prop:gtz}
The loads
\[
  L_Z(P)=\sum_{e\ni P}w_e,
  \qquad
  L_X(x)=\sum_{e\ni x}w_e,
  \qquad
  L_Y(y)=\sum_{e\ni y}w_e
\]
satisfy
\[
  \sum_{P\in Z_n}(L_Z(P)-1)^2=o(|Z_n|),                    \tag{17}
\]
\[
  \sum_{x\in X_n}(L_X(x)-L_X^{\rm lim}(x))^2=o(|X_n|),      \tag{18}
\]
and
\[
  \sum_{y\in Y_n}(L_Y(y)-L_Y^{\rm lim}(y))^2=o(|Y_n|),      \tag{19}
\]
where $L_X^{\rm lim},L_Y^{\rm lim}\le1-2\gamma$.
\end{proposition}

\begin{proof}
We use the finite-complexity Green--Tao theorem for affine-linear forms in
primes, made unconditional by the Green--Tao nilsequence theorem and the
Green--Tao--Ziegler inverse theorem
\cite{GT2010,GT2012MN,GTZ2012}.  It is cleanest here to use the normalized
auxiliary $W$-tricked form, because then the main terms are normalized residue
cell averages and no singular-series cancellation has to be checked separately.

Concretely, put $W_0=2W$.  For an auxiliary parameter $w$ let
\[
  W_{\rm aux}=W_0\prod_{\substack{\ell\le w\\ \ell\nmid W_0}}\ell .
\]
All congruence conditions modulo the fixed $W_0$ are kept as residue data, and
one then sums over residue lifts modulo $W_{\rm aux}$ on which all displayed
linear forms are units.  For each fixed $w$, the GTZ theorem gives the expected
normalized average for every fixed finite-complexity system on every rational
polytope.  After taking $n\to\infty$, one lets $w\to\infty$.  Equivalently, in
raw prime-indicator notation this says that, for a fixed finite-complexity
system $\Psi=(\psi_1,\ldots,\psi_t)$ on a fixed affine lattice and a bounded
Riemann-integrable weight $F$,
\[
  \sum_{\mathbf m\in nK\cap\Lambda}
  F(\mathbf m/n)\prod_{i=1}^t\one_{\psi_i(\mathbf m)\in\Pset}
  =\left(\mathfrak S_{\Psi,\Lambda}\int_KF+o(1)\right)
  {n^d\over(\log n)^t}.                                  \tag{20}
\]
The normalized formulation is what is used below; (20) is only a translation.
In the auxiliary residue expansion one sums only over lifts on which every
displayed form is a unit modulo $W_{\rm aux}$; the fixed typed cells modulo
$W_0$ are the residue-compatible cells already shown admissible in the
finite-core construction.

The edge-total system for a type $\tau=(a,b,\sigma,r,r',\pi)$ consists of
\[
  q,
  \qquad q',
  \qquad P_\tau(q,q')=\sigma(bq'-aq).                    \tag{21}
\]
The coefficient vectors $(1,0),(0,1),(-\sigma a,\sigma b)$ are pairwise
non-proportional.  Applying GTZ with weight $g_\tau(q/n,q'/n)$ on the support
polygon gives the first moments.  Because the typed kernels were defined by
dividing the aggregate edge intensity by the one-edge local density, these first
moments are exactly the limiting load integrals in (15c)--(15d), up to
$o(n/\log n)$.

For the label second moment, two types
\[
  \tau_i=(a_i,b_i,\sigma_i,r_i,r_i',\pi)
\]
sharing the same label are parametrized on the fixed affine lattice
\[
  a_iq_i+\sigma_iP\equiv0\pmod {b_i}
\]
by variables $(P,q_1,q_2)$, together with the full typed residue constraints
\[
  P\equiv\pi\pmod W,\qquad
  q_i\equiv r_i\pmod W,\qquad
  {a_iq_i+\sigma_iP\over b_i}\equiv r_i'\pmod W.
\]
The last congruence is part of the affine lattice; it is not implied by the
divisibility congruence when $b_i$ has factors in $W$.  The prime forms are
\[
  P,
  \quad q_1,
  \quad {a_1q_1+\sigma_1P\over b_1},
  \quad q_2,
  \quad {a_2q_2+\sigma_2P\over b_2}.                    \tag{22}
\]
Restricted to that lattice, these are integer-valued affine-linear forms.  The
coefficient vectors in $(P,q_1,q_2)$ are
\[
  (1,0,0),\ (0,1,0),\ (\sigma_1/b_1,a_1/b_1,0),\
  (0,0,1),\ (\sigma_2/b_2,0,a_2/b_2),
\]
and no two are rational multiples.

For the $X$-side second moment, two edges sharing $x=2aq$ necessarily have the
same $a$ and the same prime $q$; otherwise $a_1q_1=a_2q_2$ is impossible for
large $n$, since $a_i$ are $S$-smooth while $q_i\notin S$ are primes.  The prime
forms in variables $(q,q_1',q_2')$ are
\[
  q,
  \quad q_1',
  \quad q_2',
  \quad \sigma_1(b_1q_1'-aq),
  \quad \sigma_2(b_2q_2'-aq).                            \tag{23}
\]
The coefficient vectors are
\[
  (1,0,0),\ (0,1,0),\ (0,0,1),\
  (-\sigma_1a,\sigma_1b_1,0),\
  (-\sigma_2a,0,\sigma_2b_2),
\]
again pairwise non-proportional.  The $Y$-side moment is symmetric, with
variables $(q',q_1,q_2)$ and forms
\[
  q',
  \quad q_1,
  \quad q_2,
  \quad \sigma_1(bq'-a_1q_1),
  \quad \sigma_2(bq'-a_2q_2).                            \tag{24}
\]
These also have pairwise non-proportional coefficient vectors after the same
trivial check.

The full affine lattices in (22)--(24) include all fixed congruences modulo
$2W$ and the divisibility congruences needed to make the displayed rational
forms integral.  They are locally admissible.  At primes dividing $2W$, the
prescribed residue classes make every displayed prime form a unit.  At primes
outside $2W$, none of the coefficients $a_i,b_i$ vanish, and for each fixed
shared variable the remaining representation variables avoid at most two
residue classes.  Hence there is always an allowed residue tuple.

The only collisions excluded above are lower-dimensional repeated-edge
diagonals.  Their total weighted contribution is bounded by
\[
  O\left({\log^4 n\over n^2}\cdot {n^2\over\log^3 n}\right)
  =O(\log n)=o(n/\log n).                                \tag{25}
\]
Boundary layers are handled by first replacing every support polygon by a
slightly shrunken finite union of rational polytopes, applying GTZ there, and
then letting the shrinkage tend to zero.  Bounded piecewise-continuous kernels
are treated by $L^1$ approximation by step functions on rational polytopes.
Since the coefficient core and the type set are finite, all such errors remain
$o(n/\log n)$ after summing over types and type-pairs.

It remains to spell out why the second moments are the $L^2$ norms of the
limiting loads.  In the normalized $W_{\rm aux}$-tricked formulation, the
one-edge density assigned to a typed cell is precisely the density divided out
in the definition of $g_\tau$.  Therefore the five-form count (22), after
disintegration by $t=P/n$ and $\pi=P\bmod W$, has main term
\[
  {n\over\log n}
  \sum_{\pi\in\B}\zeta_\pi
  \int_{1/5}^{\beta}
  L_{Z,\tau_1}^{\rm lim}(t,\pi)
  L_{Z,\tau_2}^{\rm lim}(t,\pi)\,dt
  +o(n/\log n),
\]
and summing over $\tau_1,\tau_2$ gives
\[
  \sum_{P\in Z_n}L_Z(P)^2
  =
  {n\over\log n}
  \sum_{\pi\in\B}\zeta_\pi
  \int_{1/5}^{\beta}(L_Z^{\rm lim}(t,\pi))^2\,dt
  +o(|Z_n|).
\]
Since $L_Z^{\rm lim}=1$ and the prime number theorem in the fixed robust
classes gives
\[
  |Z_n|=
  {n\over\log n}\sum_{\pi\in\B}\zeta_\pi(\beta-1/5)+o(n/\log n),
\]
this proves the label second moment.  The corresponding first moment is obtained
from (21), hence (17) follows.

The $X$-side is identical in structure.  GTZ applied to (23) gives
\[
  \sum_{x\in X_n}L_X(x)^2
  =
  {n\over\log n}
  \sum_{(a,r)}\xi_{a,r}
  \int (L_X^{\rm lim}(z,a,r))^2\,dz
  +o(|X_n|).
\]
The edge-total estimate (21), now with the bounded test function
$L_X^{\rm lim}(2aQ,a,r)$ inserted, gives the mixed term
\[
  \sum_{x\in X_n}L_X(x)L_X^{\rm lim}(x)
  =
  {n\over\log n}
  \sum_{(a,r)}\xi_{a,r}
  \int (L_X^{\rm lim}(z,a,r))^2\,dz
  +o(|X_n|),
\]
and the fixed-modulus prime number theorem gives the same expression for
$\sum_x (L_X^{\rm lim}(x))^2$.  Expanding the square proves (18).  The proof of
(19) from (24) is symmetric.

If one insists on writing everything in the raw singular-series form (20), the
same argument requires checking the finite list of conditional local-factor
identities
\[
  \lambda^Z_{\tau_1,\tau_2}={\lambda_{\tau_1}\lambda_{\tau_2}\over \zeta_\pi},
  \qquad
  \lambda^X_{\tau_1,\tau_2}={\lambda_{\tau_1}\lambda_{\tau_2}\over \xi_{a,r}},
  \qquad
  \lambda^Y_{\tau_1,\tau_2}={\lambda_{\tau_1}\lambda_{\tau_2}\over \eta_{b,r'}}. \tag{26}
\]
These identities say only that, after conditioning on the shared label or side
vertex residue cell, the two remaining representation variables are locally
independent.  For primes $\ell\nmid2W$, this can be checked directly.  The
one-edge system has
\[
  (\ell-1)(\ell-2)
\]
good pairs $(q,q')\pmod\ell$, giving local factor
$\ell(\ell-2)/(\ell-1)^2$.  A label-pair system has, for each nonzero
$P\pmod\ell$, two independent choices $q_i$ each avoiding $0$ and one further
nonzero class, hence $(\ell-1)(\ell-2)^2$ good triples and local factor
$[\ell(\ell-2)/(\ell-1)^2]^2$.  The $X$- and $Y$-pair systems are identical
after conditioning on the shared side vertex.  The finitely many primes
dividing $2W$ are fixed residue-cell checks built into the affine lattices.
Thus (26) holds, and the normalized formulation above is simply a concise way
to organize the same finite verification.
\end{proof}

\section{Fractional preprocessing and Kahn rounding}

The following deterministic lemma converts Proposition~\ref{prop:gtz} into a
large fractional matching.

\begin{lemma}[Fractional preprocessing]\label{lem:preprocess}
There are subweights $0\le t_e\le w_e$ such that
\[
  \sum_{e\ni v}t_e\le1\quad\text{for every vertex }v,
  \qquad
  \sum_e t_e=(1-o(1))|Z_n|,                              \tag{27}
\]
\[
  \max_e t_e=o(1),
  \qquad
  a(t):=\max_{u\ne v}\sum_{e\supset\{u,v\}}t_e=o(1).      \tag{28}
\]
\end{lemma}

\begin{proof}
Normalize labels.  For $P\in Z_n$, set
\[
  c_P=\min(1,L_Z(P)^{-1})
\]
with $c_P=1$ if $L_Z(P)=0$, and put $w'_e=c_Pw_e$ for edges with label $P$.
Then
\[
  \sum_e w'_e=\sum_{P\in Z_n}\min(L_Z(P),1)=|Z_n|-o(|Z_n|),              \tag{29}
\]
because $1-\min(u,1)\le |1-u|$ and Cauchy--Schwarz applies to (17).
Moreover every normalized label load is at most $1$.

Since $w'_e\le w_e$, side loads only decrease.  Let
$B_X=\{x:L'_X(x)>1\}$.  Because $L_X^{\rm lim}\le1-2\gamma$, (18) implies not
only $|B_X|=o(|X_n|)$ but also
\[
  \sum_{x\in B_X}L'_X(x)=o(|Z_n|).                       \tag{30}
\]
Indeed $B_X\subset\{x:L_X(x)>1\}$, and on this latter set
$L_X(x)-L_X^{\rm lim}(x)>2\gamma$, so (18) gives
$|B_X|=o(|X_n|)$.  Also, writing
$\Delta_X(x)=L_X(x)-L_X^{\rm lim}(x)$ and using $L'_X\le L_X$,
\[
  \sum_{x\in B_X}L'_X(x)
  \le |B_X|+\sum_{x\in B_X}|\Delta_X(x)|
  \le |B_X|+|B_X|^{1/2}
  \left(\sum_{x\in X_n}\Delta_X(x)^2\right)^{1/2}
  =o(|Z_n|),
\]
since $|X_n|\asymp |Z_n|$.  The same argument gives
\[
  \sum_{y\in B_Y}L'_Y(y)=o(|Z_n|).                       \tag{31}
\]
Delete all edges incident to $B_X\cup B_Y$ and call the remaining weights
$t_e$.  Then all vertex loads are at most $1$, and (29)--(31) give
$\sum_e t_e=(1-o(1))|Z_n|$.  Also
\[
  t_e\le w_e\ll {\log^2 n\over n}=o(1).
\]
Finally the actual hypergraph has pair codegree at most $2$: a pair $(x,y)$
determines $P=|x-y|/2$, and a pair $(x,P)$ or $(y,P)$ has at most the two
opposite-side possibilities $x\pm2P$ or $y\pm2P$.  Hence
\[
  a(t)\le2\max_e t_e=o(1).
\]
\end{proof}

We now use Kahn's theorem.  In the notation of \cite{Kahn1996}, for a hypergraph
$\mathcal H$ and a fractional matching $t$, set
\[
  a(t)=\max_{x\ne y}\sum_{A\ni x,y}t(A).                  \tag{32}
\]
Kahn's Theorem 1.5 says: for fixed $k$ and $l$, if $\mathcal H$ is $k$-bounded,
$t$ is a fractional matching, and $c_1,\ldots,c_l:\mathcal H\to\mathbb R_+$
satisfy
\[
  \sum_{A\in\mathcal H}c_i(A)^2t(A)
  =o\left(\left(\sum_{A\in\mathcal H}c_i(A)t(A)\right)^2\right),        \tag{33}
\]
then as $a(t)\to0$ there is a matching $M$ with
\[
  \sum_{A\in M}c_i(A)\sim\sum_{A\in\mathcal H}c_i(A)t(A)
\]
for every $i$.

Apply this with $k=3$, $l=1$, and $c_1(e)\equiv1$.  Then
\[
  \sum_e c_1(e)^2t_e=\sum_e t_e=o\left((\sum_e t_e)^2\right)
\]
because $\sum_e t_e\asymp |Z_n|\to\infty$.  Thus there is a genuine matching
$M_n$ with
\[
  |M_n|=(1-o(1))|Z_n|.                                   \tag{34}
\]

\section{Cleanup and completion}

For each matched edge $(x,y,P)$ choose
\[
  a_P\equiv x\pmod P.
\]
Since $|x-y|=2P$, this also hits $y$.  The residue is nonzero: if $P>n/5$
divided an even residual target $x\le n$, then $x=2P$ or $x=4P$, but these are
not residual because robustness gives two $S$-hits.

Let
\[
  N=|A_S(n)|=(1+o(1)){n\over\log n}.
\]
By (9) and (34),
\[
  |M_n|=\left((\beta-1/5)\delta_S+o(1)\right)N.           \tag{35}
\]
The full robust reservoir satisfies
\[
  |\R(n)|=\left({4\over5}\delta_S+o(1)\right)N.           \tag{36}
\]
After the pair stage, the number of unused robust primes is
$|\R(n)|-|M_n|$.  The number of residual tokens still uncovered, including the
coefficient tails and all exceptional tokens, is at most $N-2|M_n|+o(N)$.
It is enough that
\[
  |\R(n)|-|M_n|\ge N-2|M_n|+o(N),
\]
or
\[
  |M_n|\ge N-|\R(n)|+o(N).                               \tag{37}
\]
This follows from the strict inequality
\[
  (\beta-1/5)\delta_S>1-{4\over5}\delta_S,
\]
which is exactly (7).

Assign each remaining residual token injectively to an unused robust prime
$P$ and choose
\[
  a_P\equiv x\pmod P.
\]
If one integer has two residual tokens, use two distinct unused robust primes.
The residue is nonzero by the same $2P,4P$ argument; odd exceptional tokens
cannot be $P$ or $3P$, because $P$ is already covered by parity plus an $S$-hit
and $3P$ is covered by parity plus the zero class modulo $3$.  The robust
side-debt lemma ensures that switching these primes creates no new unresolved
debt.

All residual tokens are defined after the fixed choices for $2$, $3$, and
$S$.  The matching and singleton assignments discharge these tokens
injectively.  The only new debts created by the cleanup primes are multiples of
those cleanup primes, and every such multiple is handled by the robust side-debt
lemma.  All other primes remain in their zero classes, $2$ remains in the odd
class, $3$ remains in the zero class, and the primes in $S$ remain in their
fixed nonzero classes.  The exact debt inequality (1) is satisfied for every
$m\le n$, so every $m\in[1,n]$ lies in at least two chosen residue classes.

\begin{theorem}
For all sufficiently large $n$, there are residue classes $a_p\pmod p$ for all
primes $p\le n$ such that every integer $m\in[1,n]$ satisfies
\[
  \#\{p\le n:m\equiv a_p\pmod p\}\ge2.
\]
\end{theorem}

\begin{thebibliography}{9}

\bibitem{GT2010}
B. Green and T. Tao,
\emph{Linear equations in primes},
Annals of Mathematics \textbf{171} (2010), 1753--1850.

\bibitem{GT2012MN}
B. Green and T. Tao,
\emph{The M\"obius function is strongly orthogonal to nilsequences},
Annals of Mathematics \textbf{175} (2012), 541--566.

\bibitem{GTZ2012}
B. Green, T. Tao, and T. Ziegler,
\emph{An inverse theorem for the Gowers $U^{s+1}[N]$-norm},
Annals of Mathematics \textbf{176} (2012), 1231--1372.

\bibitem{Kahn1996}
J. Kahn,
\emph{A linear programming perspective on the Frankl--R\"odl--Pippenger theorem},
Random Structures \& Algorithms \textbf{8} (1996), 149--157.

\end{thebibliography}

\end{document}
